% Shared preamble for the Quantum Computing companion lecture notes.
\usepackage[margin=1in]{geometry}
\usepackage{amsmath,amssymb,amsthm,mathtools}
\usepackage{braket}          % \ket, \bra, \braket, \ketbra
\usepackage{enumitem}
\usepackage{graphicx}
\usepackage{xcolor}
\usepackage[colorlinks=true,linkcolor=blue,citecolor=blue,urlcolor=blue]{hyperref}

% ---- Theorem-like environments (shared counter within each chapter) ----
\theoremstyle{definition}
\newtheorem{definition}{Definition}[chapter]
\newtheorem{example}[definition]{Example}
\theoremstyle{plain}
\newtheorem{theorem}[definition]{Theorem}
\newtheorem{lemma}[definition]{Lemma}
\newtheorem{proposition}[definition]{Proposition}
\newtheorem{corollary}[definition]{Corollary}
\theoremstyle{remark}
\newtheorem{remark}[definition]{Remark}

% Exercises use their own per-chapter counter.
\theoremstyle{definition}
\newtheorem{exercise}{Exercise}[chapter]

% Solution environment (unnumbered, italic "Solution." lead-in).
\newenvironment{solution}{\begin{proof}[\normalfont\itshape Solution.]}{\end{proof}}

% ---- Common macros ----
\newcommand{\C}{\mathbb{C}}
\newcommand{\R}{\mathbb{R}}
\newcommand{\Z}{\mathbb{Z}}
\newcommand{\N}{\mathbb{N}}
\newcommand{\F}{\mathbb{F}}
\DeclarePairedDelimiter{\abs}{\lvert}{\rvert}
\DeclarePairedDelimiter{\norm}{\lVert}{\rVert}
\newcommand{\Tr}{\operatorname{Tr}}
\newcommand{\id}{\mathbb{1}}
\newcommand{\CNOT}{\mathrm{CNOT}}
% braket.sty ships \ketbra commented out; define it here so all chapters have it.
\providecommand{\ketbra}[2]{\ket{#1}\!\bra{#2}}
\newcommand{\Had}{H}
\newcommand{\QFT}{\mathrm{QFT}}
\newcommand{\Uf}{U_f}
\newcommand{\defeq}{\vcentcolon=}
